\uuid{8RPF}
\titre{Composée/DL usuels}
\auteur{hanani ; mzari}
\organisation{exo7}
\licence{CC BY-NC-SA 4.0}
\datecreate{2020-01-20}
\niveau{L1}
\module{Analyse}
\chapitre{Développement limité}
\sousChapitre{Calculs}
\difficulte{2}
\theme{Opérations sur les DL, Composée, DL usuels}
\exo7qcmid{lille-2.developpements-limites.009}

\contenu{
\texte{
Soit $a\in \Rr$ et $f$ une fonction. On dit que $f(x)$ est un $o\left((x-a)^n\right)$ au voisinage de $a$ et on écrit $f(x)=o((x-a)^n)$ si $\displaystyle \lim _{x\to a}\frac{f(x)}{(x-a)^n}=0$. On note par $DL_n(a)f(x)$ le développement limité de $f$ à l'ordre $n$ en $a$. On note aussi $DL_n(a^+)f(x)$ (resp. $DL_n(a^-)f(x)$) le développement limité de $f$ à l'ordre $n$ à droite en $a$ (resp. à gauche en $a$).
}

\begin{question}
\qclassification{125.02}
\qtags{motcle=Composée/DL usuels}
Les égalités suivantes portent sur des développements limités en $0$. Cocher celles qui sont vraies :
\begin{answers}
\bad{$\displaystyle \ln(1+2x)=2\left(x-\frac{x^2}{2}+o(x^2)\right)$.}
\good{$\displaystyle \ln(1-2x)=-2x-2x^2+o(x^2)$.}
\good{$\displaystyle \ln(1+x^2)=x^2-\frac{x^4}{2}+o(x^4)$.}
\bad{$\displaystyle \ln(1-x^2)=-\left[x^2-\frac{x^4}{2}+o(x^4)\right]$.}
\end{answers}
\begin{explanations}
On applique le développement : $\displaystyle DL_2(0)\ln(1+t)=t-\frac{t^2}{2}+o(t^2)$ avec $t=2x$, $t=-2x$, $t=x^2$ et enfin avec $t=-x^2$.
\end{explanations}
\end{question}
}
